\documentclass[
    language=english,
    doctype=postdoc-report,
    institution=none,
]{../../omnilatex}

\RequirePackage{config/document-settings}

\begin{document}
\maketitle

\begin{abstract}
This annual report summarizes research activities, publications, teaching
contributions, and professional development during the second year of
my postdoctoral fellowship at the Center for Computational Biology,
Department of Computer Science, University of Exampleton. The primary
research focus has been on developing scalable Bayesian methods for
spatial transcriptomics data analysis, with applications to the human
brain atlasing project. Key achievements include 4 peer-reviewed
publications (2 as first author), 2 manuscripts under review, and
the release of an open-source software package with over 500 GitHub
stars. Teaching contributions include co-instructing a graduate course
on probabilistic machine learning and mentoring 3 PhD students.
\end{abstract}

\section{Research Summary}
\label{sec:research}

\subsection{Project 1: Bayesian Spatial Transcriptomics (BST)}

\subsubsection{Background and Motivation}
Spatial transcriptomics technologies (MERFISH, Slide-seq, Visium) measure
gene expression while preserving spatial context in tissue sections. These
technologies generate datasets with millions of spatially resolved
measurements, presenting significant computational challenges for
probabilistic modeling \cite{stahl2016visualization}.

\subsubsection{Objectives}
\begin{enumerate}
    \item Develop scalable variational inference algorithms for spatial
          gene expression models that scale to datasets with $>$ 1M spots.
    \item Implement a Gaussian process latent variable model (GPLVM) that
          captures spatially varying gene expression patterns.
    \item Validate the methods on the Human Brain Atlas consortium dataset
          (120 brain sections, 50,000--200,000 spots per section).
\end{enumerate}

\subsubsection{Methods and Results}

I developed BST-spGPLVM, a Bayesian spatial Gaussian process latent variable
model that uses inducing point approximations and natural gradient variational
inference to achieve $\mathcal{O}(nm)$ complexity, where $n$ is the number
of spots and $m$ is the number of inducing points.

\begin{table}[htbp]
    \centering
    \caption{Scalability comparison of BST-spGPLVM against existing methods.}
    \label{tab:scalability}
    \begin{tabular}{@{}lcccc@{}}
        \toprule
        Method & Spots & Genes & Time (min) & Memory (GB) \\
        \midrule
        SpatialDM & 10K & 20K & 12.4 & 8.2 \\
        SPARK-X & 10K & 20K & 45.6 & 12.1 \\
        BayesSpace & 10K & 20K & 28.3 & 9.8 \\
        BST-spGPLVM & 10K & 20K & \textbf{3.2} & \textbf{2.1} \\
        \midrule
        BST-spGPLVM & 100K & 20K & \textbf{28.7} & \textbf{8.4} \\
        BST-spGPLVM & 1M & 20K & \textbf{245} & \textbf{42.3} \\
        \bottomrule
    \end{tabular}
\end{table}

The model identifies spatially variable genes (SVGs) with higher precision
than existing methods (AUC = 0.92 vs.\ 0.78 for SpatialDM on simulated
data with known SVGs). Application to the Human Brain Atlas data reveals
previously uncharacterized spatial gene expression gradients in the
prefrontal cortex.

\subsubsection{Publications}

\begin{enumerate}
    \item \textbf{[First author]} ``Scalable Bayesian Spatial Transcriptomics
          with Gaussian Process Latent Variable Models,'' \textit{Nature
          Methods}, 2025. (Under review)
    \item \textbf{[Co-author]} ``Bayesian Methods for Spatial Genomics:
          A Comparative Benchmark,'' \textit{Bioinformatics}, 2025.
\end{enumerate}

\subsection{Project 2: Hierarchical Bayesian Models for Brain Atlasing}

\subsubsection{Background}
The Brain Atlas Consortium is generating spatial transcriptomics data from
1,000+ postmortem brain samples to create a comprehensive reference atlas
of human brain gene expression.

\subsubsection{Objectives}
\begin{enumerate}
    \item Develop hierarchical Bayesian models that borrow strength across
          brain regions and donors to improve estimation accuracy.
    \item Design an active learning framework to optimally select which
          brain regions to sample next.
    \item Create a data integration pipeline that combines spatial
          transcriptomics with single-cell RNA-seq and proteomics data.
\end{enumerate}

\subsubsection{Methods and Results}

I developed a hierarchical Gaussian process model that shares information
across brain regions through a common covariance structure. The model
uses stochastic variational inference with distributed computing to
handle the full atlas dataset.

\begin{equation}
    f_s(\mathbf{x}) \sim \mathcal{GP}\left(0, \kappa_s(\mathbf{x}, \mathbf{x}') + \kappa_{\text{shared}}(\mathbf{x}, \mathbf{x}')\right)
    \label{eq:hierarchical}
\end{equation}

where $\kappa_s$ is a region-specific kernel and $\kappa_{\text{shared}}$
is a shared kernel across regions. The hierarchical structure reduces
posterior variance by 23\% compared to region-independent models, as
measured by cross-validation log-likelihood.

\subsubsection{Publications}

\begin{enumerate}
    \item \textbf{[First author]} ``Hierarchical Bayesian Modeling of
          Spatial Gene Expression Across the Human Brain,'' \textit{Nature
          Neuroscience}, 2025. (Submitted)
    \item \textbf{[Co-author]} ``Active Learning for Spatial Transcriptomics
          Experimental Design,'' \textit{Genome Biology}, 2025.
\end{enumerate}

\subsection{Project 3: Open-Source Software Development}

I developed and maintain the \texttt{SpatialBayes.jl} package, a Julia
library for Bayesian spatial transcriptomics analysis. Key features:

\begin{itemize}
    \item Scalable Gaussian process inference with inducing points
    \item Hierarchical modeling across multiple tissue sections
    \item GPU acceleration via CUDA.jl
    \item Interactive visualization dashboards
    \item Python and R bindings via PyCall and RCall
\end{itemize}

\begin{table}[htbp]
    \centering
    \caption{Software metrics as of December 2025.}
    \label{tab:software}
    \begin{tabular}{@{}lc@{}}
        \toprule
        Metric & Value \\
        \midrule
        GitHub stars & 523 \\
        Forks & 87 \\
        Contributors & 12 \\
        Downloads (JuliaHub) & 8,420 \\
        Test coverage & 94\% \\
        Documentation pages & 156 \\
        Issues closed & 43 \\
        \bottomrule
    \end{tabular}
\end{table}

\section{Publications}
\label{sec:publications}

\subsection{Journal Articles}

\begin{enumerate}
    \item \textbf{[First author]} [Author], [Co-author A], [Co-author B].
          ``Scalable Bayesian Spatial Transcriptomics with Gaussian Process
          Latent Variable Models.'' \textit{Nature Methods} (under review,
          2025).

    \item \textbf{[First author]} [Author], [Co-author A]. ``Hierarchical
          Bayesian Modeling of Spatial Gene Expression Across the Human
          Brain.'' \textit{Nature Neuroscience} (submitted, 2025).

    \item [Co-author C], \textbf{[Author]}, [Co-author D]. ``Bayesian
          Methods for Spatial Genomics: A Comparative Benchmark.''
          \textit{Bioinformatics} 41(3): btae012, 2025.

    \item [Co-author E], \textbf{[Author]}, [Co-author F]. ``Active
          Learning for Spatial Transcriptomics Experimental Design.''
          \textit{Genome Biology} 26: 45, 2025.
\end{enumerate}

\subsection{Conference Papers}

\begin{enumerate}
    \item \textbf{[First author]} [Author], [Co-author A]. ``BST-spGPLVM:
          Scalable Bayesian Spatial Transcriptomics.'' \textit{Proceedings
          of NeurIPS 2025} (submitted).

    \item \textbf{[First author]} [Author]. ``Distributed Variational
          Inference for Large-Scale Spatial Genomics.'' \textit{Proceedings
          of AISTATS 2026} (submitted).
\end{enumerate}

\subsection{Preprints}

\begin{enumerate}
    \item \textbf{[First author]} [Author], [Co-author A], [Co-author B].
          ``SpatialBayes.jl: A Julia Library for Bayesian Spatial
          Transcriptomics Analysis.'' arXiv:2025.12345, 2025.
\end{enumerate}

\section{Teaching and Mentoring}
\label{sec:teaching}

\subsection{Course Instruction}

\begin{table}[htbp]
    \centering
    \caption{Teaching contributions during the reporting period.}
    \label{tab:teaching}
    \begin{tabular}{@{}lcccc@{}}
        \toprule
        Course & Role & Semester & Enrollment & Rating \\
        \midrule
        CS 7641: Probabilistic ML & Co-instructor & Fall 2025 & 38 & 4.7/5.0 \\
        CS 6250: Bioinformatics & Guest lectures (3) & Spring 2025 & 62 & 4.8/5.0 \\
        CS 8903: Bayesian Methods & Seminar co-leader & Fall 2025 & 12 & --- \\
        \bottomrule
    \end{tabular}
\end{table}

\subsection{Student Mentoring}

\begin{enumerate}
    \item \textbf{Alex Kim} (PhD student, Year 2): Co-mentoring on spatial
          transcriptomics methods. Alex presented a poster at ISMB 2025
          and is preparing a first-author manuscript.

    \item \textbf{Priya Sharma} (PhD student, Year 1): Mentoring on
          variational inference methods. Priya developed a novel divergence
          measure for spatial data, accepted at AISTATS 2026.

    \item \textbf{Marcus Lee} (MS student, final year): Thesis supervision
          on GPU-accelerated Gaussian process inference. Thesis defended
          August 2025.
\end{enumerate}

\subsection{Undergraduate Research}

I supervised 2 undergraduate research projects through the university's
REU (Research Experiences for Undergraduates) program:

\begin{itemize}
    \item \textbf{Jordan Rivera} (Summer 2025): Developed a visualization
          toolkit for spatial gene expression data. Project resulted in
          a poster presentation at the undergraduate research symposium.
    \item \textbf{Casey Wong} (Fall 2025): Benchmarking spatial
          transcriptomics methods on simulated data. Ongoing.
\end{itemize}

\section{Professional Development}
\label{sec:development}

\subsection{Conference Participation}

\begin{itemize}
    \item NeurIPS 2025 (Vancouver, Canada): Poster presentation
    \item ISMB 2025 (Vienna, Austria): Oral presentation (selected)
    \item Single-Cell Biology 2025 (Cambridge, UK): Invited speaker
    \item Bayesian Statistics in Science (BSiS) 2025: Panel member
\end{itemize}

\subsection{Workshops and Tutorials}

\begin{itemize}
    \item Co-organized ``Bayesian Methods for Spatial Genomics'' workshop
          at NeurIPS 2025 (45 attendees)
    \item Presented tutorial on ``Gaussian Processes for Spatial Data''
          at ISMB 2025 (120 attendees)
    \item Guest lecture at University of Cambridge, Department of Statistics
\end{itemize}

\subsection{Service}

\begin{itemize}
    \item Program committee: NeurIPS 2025, ICML 2026
    \item Reviewer: Nature Methods, Bioinformatics, JMLR
    \item Departmental seminar organizer (2025--2026 academic year)
    \item Faculty search committee member (junior faculty position)
\end{itemize}

\section{Future Plans}
\label{sec:future}

\subsection{Research Directions}

\begin{enumerate}
    \item \textbf{Foundation models for spatial biology:} Develop a
          pre-trained model for spatial transcriptomics data, analogous
          to protein language models (ESM-2, AlphaFold), that learns
          general representations of spatial gene expression patterns.

    \item \textbf{Multi-modal spatial integration:} Extend BST-spGPLVM
          to jointly model spatial transcriptomics with spatial proteomics
          (CODEX, MIBI) and spatial epigenomics (Spatial-CUT&Tag).

    \item \textbf{Clinical applications:} Apply the methods to spatial
          transcriptomics data from brain tumor biopsies, in collaboration
          with the Neurosurgery Department, to develop diagnostic and
          prognostic biomarkers.
\end{enumerate}

\subsection{Career Goals}

\begin{table}[htbp]
    \centering
    \caption{Career development milestones and timeline.}
    \label{tab:career}
    \begin{tabular}{@{}ld{12}l@{}}
        \toprule
        Timeline & Milestone \\
        \midrule
        Q1 2026 & Submit Nature Methods first-author manuscript \\
        Q2 2026 & Submit NIH R21 grant (spatial biology methods) \\
        Q3 2026 & Apply for faculty positions \\
        Q4 2026 & Finalize tenure-track applications \\
        2027    & Begin tenure-track position (target) \\
        \bottomrule
    \end{tabular}
\end{table}

\subsection{Job Market Preparation}

I am preparing for the academic job market in Fall 2026. Key preparation
activities include:

\begin{itemize}
    \item Developing a 3-page research statement focused on ``Bayesian
          Methods for Spatial Biology''
    \item Preparing a teaching portfolio with syllabus designs for 3
          proposed courses
    \item Networking with department chairs at target institutions
\end{itemize}

\section{Summary of Achievements}
\label{sec:summary}

\begin{table}[htbp]
    \centering
    \caption{Summary of key metrics for the reporting period.}
    \label{tab:summary}
    \begin{tabular}{@{}lc@{}}
        \toprule
        Category & Metric \\
        \midrule
        Publications & 4 published, 2 under review \\
        First-author papers & 2 \\
        Software (GitHub stars) & 523 \\
        Teaching & 1 course, 3 guest lectures \\
        Students mentored & 5 (3 PhD, 1 MS, 1 REU) \\
        Invited talks & 3 \\
        Conference presentations & 4 \\
        Grant applications & 1 submitted (NIH R01) \\
        \bottomrule
    \end{tabular}
\end{table}

This has been a productive year. The core research on scalable Bayesian
spatial transcriptomics has advanced significantly, with two first-author
manuscripts under review at top journals. The software package has gained
traction in the community, and teaching and mentoring activities have been
rewarding. I look forward to building on these achievements in the coming
year, with a focus on foundation models for spatial biology and preparation
for the faculty job market.

\bibliographystyle{plain}
\bibliography{references}

\end{document}
